%% kapa: exam=1 round=30 year=2019 period=2 form=A subject=appraisal_law subject_ko=감정평가관계법규 seq=4 number=4
%% source: Q-Net 감정평가사 기출문제 (cst00309) artlSeq=5205205 file=extracted/1008/5205205/2019년 제30회 감정평가사 1차시험 문제지/2019년 제30회 감정평가사 1차시험 문제지/2019년도 제30회 감정평가사 1차 2교시 A형.hwp
%% conversion: kordoc 4.24.51 -> tools/gen1.py (auto); review=proofread-sonnet; flags=none
\begin{question}{4}
국토의 계획 및 이용에 관한 법령상 용도지역별 건폐율의 최대한도가 큰 순서대로 나열된 것은? (단, 조례 및 기타 강화ㆍ완화조건은 고려하지 않음)
\begin{framedtext}
\begin{tabularx}{\linewidth}{@{}X@{\hspace{1em}}X@{}}
ㄱ. 계획관리지역 & ㄴ. 자연녹지지역 \\
ㄷ. 근린상업지역 & ㄹ. 제2종일반주거지역
\end{tabularx}
\end{framedtext}
\begin{choices}
  \item ㄷ-ㄹ-ㄱ-ㄴ
  \item ㄷ-ㄹ-ㄴ-ㄱ
  \item ㄹ-ㄱ-ㄷ-ㄴ
  \item ㄹ-ㄷ-ㄱ-ㄴ
  \item ㄹ-ㄷ-ㄴ-ㄱ
\end{choices}
\end{question}
\answer{1}
